\clearpage
\section{Glosario}\label{sec:spmp-7}
Los términos siguientes se utilizan con el significado operativo indicado en este plan. Las siglas de estándares identifican las fuentes citadas en las secciones donde se aplican y no se emplean como garantía de cumplimiento. Los conceptos de simulación se incluyen únicamente en la medida necesaria para comprender la administración de los productos y sus evidencias.
\begin{longtable}{|p{0.29\dimexpr\linewidth-23pt\relax}|p{0.71\dimexpr\linewidth-23pt\relax}|}
\caption{Glosario de gestión y del producto. Fuente: síntesis de los materiales internos citados en el plan.}\\
\hline
\textbf{Término} & \textbf{Significado en el proyecto}\\\hline
\endfirsthead
\hline
\textbf{Término} & \textbf{Significado en el proyecto}\\\hline
\endhead
DAD & Disciplined Agile Delivery, enfoque elegido para organizar Inicio, Construcción y Transición de forma iterativa e incremental.\\\hline
FDS & Fire Dynamics Simulator. Software de simulación empleado para preparar los escenarios de incendio.\\\hline
WFDS & Wildland-Urban Interface Fire Dynamics Simulator. Extensión histórica de FDS para incendios forestales y de interfaz urbano-forestal, cuyas capacidades se integraron en FDS.\\\hline
Smokeview & Herramienta seleccionada para presentar las animaciones georreferenciadas de los resultados FDS.\\\hline
HPC y ZINE & Cómputo de alto desempeño e infraestructura universitaria prevista para la ejecución científica.\\\hline
Biblioteca de escenarios & Conjunto catalogado de resultados precalculados con sus condiciones, procedencia y límites.\\\hline
Igniciones de referencia & Puntos de inicio de escenarios con densidad variable guiada por amenaza. No son las celdas de la malla CFD.\\\hline
Malla CFD & Discretización del dominio físico dentro de cada simulación.\\\hline
API & Interfaz prevista para obtener información meteorológica desde un servicio externo.\\\hline
Incremento & Resultado revisable del trabajo, que puede incorporar software o evidencia técnica.\\\hline
WBS & Estructura de descomposición del trabajo que relaciona productos y actividades.\\\hline
Línea base & Versión identificada que se acuerda como referencia para estudiar y controlar cambios.\\\hline
Trazabilidad & Relación entre origen, requisito, producto, versión y evidencia de comprobación.\\\hline
Verificación y validación & Comprobación de condiciones especificadas y evaluación de adecuación al uso previsto, respectivamente.\\\hline
BPMN y UML & Notación para describir procesos de gestión y lenguaje para modelar el software, con propósitos distintos.\\\hline
Cliente o interlocutor usuario & Persona con quien se acordarán necesidades y aceptación. Su identificación no presupone una institución contratante.\\\hline
\end{longtable}
